% 原书第 93 页（PDF 第 106 页）
\BookPage{93}{106}
\BookHeading{中国社会思想视野中的「难得糊涂」和「郑板桥现象」}{中国社会思想视野中的「难得糊涂」和「郑板桥现象」}{吴建华*}
\begin{BookQuote}
「难得糊涂」是郑板桥格言。这格言和响亮的「郑板桥」的名号一起，几乎成
\end{BookQuote}

为家喻户晓的中华文化常识。从社会思想的角度看，这格言和「郑板桥现象」，本身就是中国社会思想研究的极好题材。

郑燮（1693-1765），字克柔，号板桥，清代扬州府兴化县（今江苏泰州市兴化市）人，不仅是我国杰出的思想家、文学家，而且是书画艺术大家，「扬州八怪的主将。他身历清代小农「盛世」，是「康熙秀才，雍正举人，乾隆进士」，在山东范县（今河南濮阳市范县）、潍县（今山东潍坊市）当过12年县令。他为民请赈，得罪了上司大吏，被参劾罢官，解职归里。总体上，他任官时间较短，且沉沦于知县下吏，不善逢迎，两袖清风，也失去了升迁机会，一生穷困潦倒，因此，他大多数时间只得依靠出卖自己的字画谋生，是一介书生，近于平民的民间艺术家。

然而，板桥留给世人大量的书画作品，这些作品本身连同上面的题跋，和他的《家书》、《诗钞》、《词钞》、《小唱》及题跋碑记等文字，都反映出他丰富的社会阅历，愤世嫉俗的真实情感，以及率性诚挚的人品，蕴含着可贵独特的社会思想，不仅值得人们从文学艺术的角度去深入探讨，而且应当引起社会思想界的重视。

\vfill
\begin{BookNotes}
\BookNote{*}{作者任职于苏州大学社会学院。}
\BookNote{1.}{《郑板桥集》《补遗·板桥自叙》，第177页，上海古籍出版社，1962年1月第1版，1979年新1版，又见于板桥印章。关于郑板桥的文字，真伪难辨之处甚多，研究使用时候尤当注意。本文以上海古籍出版社1962年版的《郑板桥集》为准。卞孝萱先生积数十年功力，倾心寻考其故乡先哲板桥的文字，成果集中见其所编《郑板桥全集》（齐鲁书社，1985年版)、着《郑板桥丛考》（辽海出版社，2003年版；又收入《卞孝萱文集》，凤凰出版社，2010年版)。}
\end{BookNotes}

% 原书第 94 页（PDF 第 107 页）
\BookPage{94}{107}
郑板桥的社会思想至少集中体现在社会平等观、社会批判思想、社会治理思想三个方面，反映清代「盛世」时期扬州与其他地区多方面的社会现象，板桥深刻的社会观察意识、批判能力，以及他理想的社会目标、社会治理的实践活动，在中国社会思想史上具有重要的价值。2

从郑板桥的社会平等观、他对于所处的清代小农「盛世」社会现实的批判、社会治理思想、社会思想实践诸方面来看，毫无疑问，板桥的脑子是十分清醒的，充满了一腔为民请命、为民造福的士人君子的良知，反映他深受儒家入世学说的影响，修身、齐家、治国、平天下，层层递进，真要一展宏图大志。

可是一到「难得糊涂」的境地，就有了一番讲究。「难得糊涂」，一是可以理解成一直不糊涂，偶而糊涂一下，真是难得！遇着不顺心之事，也不妨开一眼，闭一眼，糊涂一次，很快就过去了，完全不必拘泥于某些原则伦理，甚至规条；二则也可以理解成要糊涂很难：因为理智的清醒，所以想要糊涂一点，很难啊！

倘若板桥的真实心态属于前一种情况，则与他勇于作为的社会思想及其实践相比，陡然转了一百八十道弯，由进取转向了退隐、逃遁。这是由儒家思想占据上风转变到让道家顺其「自然」的处世思想占据了上风。「难得糊涂」，可以成为一种掩饰，既能是失败的掩饰，也可是清醒的掩饰，正是一种良知的颠倒，社会正义的颠倒。它是一种处世哲学，一种退让，明哲保身，自我安慰，求得心理平衡。实际上，正是心理与社会难以平衡的极好外露。

然而，板桥的真实心态属于后一种情况，无疑体现了板桥一贯的为人处世风格。他常见的「难得糊涂」条幅注脚清楚地说：「聪明难，糊涂难，由聪明而转入糊涂更难。放一着，退一步，当下心安，非图后来福报也。」他寄望的是由聪明「难得糊涂」一下。但是，他根本难以做到。因为，人们日日夜夜世世代代企求的总是聪明，不是糊涂，像板桥绝顶的聪明，要以糊涂面目出现，实在太难，难于上青天了。这正是他有感而发，是他一生阅历的深刻体知。因此，「难得糊涂」，语

\vfill
\begin{BookNotes}
\BookNote{2.}{详见吴建华、张明、刘慧，《郑板桥的社会思想》，载陆学艺、王处辉主编《追寻中国社会的自性中国社会思想史论集》，广西人民出版社，2004年版。}
\BookNote{3.}{《板桥家书—糊涂成功大全》，首页图版，中国对外翻译出版公司，2001年版。}
\end{BookNotes}

% 原书第 95 页（PDF 第 108 页）
\BookPage{95}{108}
\noindent 词的模糊遮掩不了他内心对现实的极度愤懑，以及决不妥协的勇士斗志，犹如他笔下一再渲染寓于深意的弱细幽兰，虽然雨打风霜，泥巴满身，东倒西歪，则绝没有一副媚骨奴态，依旧迎冰霜雪寒，向人间飘洒着一股淡淡的幽香，启人心扉，昭示着春天的讯息终将长存天地之间。这才是真正的板桥，「怪」的板桥，「二十年前旧板桥」4！世道变幻，而板桥终于不变，也难于「糊涂」地去变。因此，「难得糊涂」成了板桥顿悟的座右铭，让他在处世上更领有伸缩度，也更容易被人们接受。而不同层次的人对于他的格言，可以出于自己的身分地位，作自己的理解，体悟人生社会的哲学，认知社会的变迁。

应该说，板桥的「怪」的高明过人之处于此，板桥的「怪」即独特的社会思想精华于此。板桥的社会平等观、社会批判观、社会治理观、社会思想的实践都渗透着他「怪」的「难得糊涂」的意识，只是有时有意识、有时无意识地表现罢了。而终究是以无意识的「难得糊涂」为上，以有意识的「难得糊涂」为下。他能精确地把握住清代社会的特征，其实可以说相当多的也是传统中国社会的特征，从而提炼出非常精辟的处世之道来一「难得糊涂」。这格言高度糅合了儒道两家的思想精髓，模糊的词语，不知道道出了多少人的心声，概括了多少社会治理的甜酸苦辣。社会治理得好，是清醒的社会实践；治得不好呢，就算是「难得糊涂」的社会实践吧。也不要苛求治理者了吧。然而，不可否定，有时也可以取得用貌似的「难得」的糊涂，无为的方式，即假装糊涂换得的社会治理实践的良效。

「难得糊涂」让人各有所需所取，选取自如，大智若愚者亦可，真智真愚者亦可，社会就如此流传着郑板桥，他的名字、他自号的怪「六分半书」、他的「难得糊涂格言等等。「郑板桥现象」成为清代以来的社会现象，当然，某种程度上，他也是郑板桥以前的与「郑板桥现象」相同的社会现象的抽理与总结。

从艺术史的角度看，「怪」，就是「扬州八怪」画派的艺术特色，当然，它也是郑板桥个性的重要特色，是他社会思想的重要特色。没有「怪」，就没有郑板桥，就没有「扬州八怪」的艺术成就和瑰宝。实际上，这种「怪」，并不是天生的奇异，刻意的做作，而是天性的喜好自然地流露在艺术创新上，而是时代正好处

\vfill
\begin{BookNotes}
\BookNote{4.}{秦岭云编着，《扬州八家丛话》，郑燮印鉴，上海人民美术出版社，1985年版，第94页。}
\end{BookNotes}

% 原书第 96 页（PDF 第 109 页）
\BookPage{96}{109}
\noindent 于呼唤艺术大师出现的当口而出现的现象，但又是社会上大多数人仍然处于睡眠的状态。因而，正统的艺术界和大多数人见到「扬州八怪」的书画作品，则目之为「怪」，就是别致，而别致，既可含有痛骂鄙夷排斥的意图，又着实令人注目，耳目一新，起到领变先驱功能。在这个意义上，我们考察郑板桥的社会思想应该具有同样的视野，所谓板桥「怪」的社会思想，实际上是见怪不怪为好。一旦社会上普遍认为板桥的社会思想见怪不怪，岂不意味着社会的变化已经深入人心，对于板桥的认同等级和状态几乎一致了，郑板桥的意识就成为社会的意识了。

诚然，郑板桥的社会思想既充满了正直的良知，希求着理想的仁义社会目标，闪烁着人性的光辉，同时不可讳言，它贯穿着小农思维，因此仍然不可避免存在着糟粕。例如，他的社会平等观和重农抑士的过激思想是对立的，也和他自身的士人地位相抵触。他很早就靠出卖字画为生，一生都依赖书画市场改善生计，包括在山东为官时期，也深知「千家养女先教曲，十里栽花算种田」的烟花扬州才是他兴化家乡所处的经济文化行政中心，那儿商贾云集，人才济济，良好的人文和物质环境更有利于他的个人成才以及艺术发展，可是，他却又在骨子和行为上轻视并调侃那些需要字画出得起价的富商和部分官宦。不过，应当看到，正是这种农业文明的思维才使板桥能够找到评析社会的最好的「旧」武器。他用它解剖社会的黑暗，抨击社会的不公，继承着我国朴素民主思想的精华。

板桥的文学分析思维、文学主张与历史环境、社会实况紧密相连，寓社会思想于文学艺术之中，也是他社会思想的又一大特色。他高度赞赏和效仿的杜甫诗，「只一开卷阅其题次，一种忧国忧民忽悲忽喜之情，以及宗庙丘墟，关山劳戍之苦，宛然在目。其题如此，其诗有不痛心入骨者乎！」反之，他憎恶腐败的朝政，同时深知陆游诗的蕴义，有特定的时代性。「南宋时，君父幽囚，栖身杭越，其辱与危亦至矣。讲理学者推极于毫厘分寸，而卒无救时济变之才；在朝诸大臣皆流连诗酒，沉溺湖山，不顾国之大计。…且国将亡，必多忌。躬行桀、纣，必曰驾尧、舜而轶汤、武。宋自绍兴以来，主和议，增岁币，送尊号，处卑朝，刮民膏，戮大

\vfill
\begin{BookNotes}
\BookNote{5.}{《郑板桥集》，〈扬州〉，第30页。}
\BookNote{6.}{参见吴建华，〈郑板桥与徽州考述〉，《清史研究》2004年第4期；〈郑板桥与江南的关系〉，《社会科学》2006年第8期。}
\end{BookNotes}

% 原书第 97 页（PDF 第 110 页）
\BookPage{97}{110}
\noindent 将，无恶不作，无陋不为。百姓莫敢言喘，放翁恶得形诸篇翰以自取戾乎！」论诗的情感同样反映他社会思想的倾向。

在传统社会世风日下的时候，板桥看到了社会阴黑，人性堕落，因而宁愿固守着儒家君子风度，与社会世风格格不入，由此反而变得「奇异」起来，成为「怪的典型了。清人郑方坤理解板桥很透，不愧为板桥的知音：「板桥徒以狂故不理于口，然其为人内行醇谨，胸中具有泾渭。所刻寄弟书数纸皆老成忠厚之言，大有光禄《庭诰》、《颜氏家训》遗意，异乎放荡以为高者，信贤者之不可测也。昔晋文王称阮嗣宗为至慎，吾于板桥亦云。」因此，板桥的所谓怪，正是他社会思想的独特之处，板桥以及「扬州八怪」具有的社会思想地位。「难得糊涂」的怪，就是一直清醒的良知备受社会煎熬之后的无奈选择，也是社会为良知的长存付出的沉重代价。社会应当为此种社会思想的产生和普遍被国人接受承担罪责，并且要及早觉醒，积极地振作起来，要么真糊涂，要么不糊涂，何必要「难得糊涂」呢！

\vfill
\begin{BookNotes}
\BookNote{7.}{《郑板桥集》，〈范县署中寄舍弟墨第五书〉，第14页。}
\BookNote{8.}{《郑板桥集》，辑《国朝耆献类征初编》卷233，郑方坤〈郑燮小传〉，第240-241页。}
\end{BookNotes}

% 原书第 98 页（PDF 第 111 页）
\BookPage{98}{111}
